\chapter*{Abstract}
\addcontentsline{toc}{chapter}{Abstract}

Corrosion of steel reinforcement is a principal durability concern for
cementitious structural elements, and repeated photography offers an
inexpensive, non-destructive way to observe it over time. This activity
uses a predecessor experimental dataset of 48 ferrocement specimens tested
across two exposure campaigns and photographed repeatedly; in the current
audited basis, 791 aligned, readable observations carry dense
visible-corrosion labels, while wire-area loss and ultimate load are each
obtained only once per specimen, at the end of its exposure period. Built
around this asymmetry between dense surface supervision and sparse
terminal structural supervision, the activity progressed through
classical, deep-embedding, and interpretable image representations of
visible corrosion; a feasibility study of hidden structural-condition
inference; robustness evaluation under grouped and transfer-based holdout
regimes; a refocused study of terminal ultimate load; and an exploratory
degradation and proxy-remaining-useful-life screening chain. Visible
surface corrosion proved the strongest and most consistently learnable
image-based quantity, although specific benchmarks carry a
label-adjacency risk and treatment-level transfer robustness proved
phase-dependent. Hidden structural inference was substantially weaker and
less transferable, and must be interpreted in light of its sparse,
terminal, campaign-confounded supervision. Image-derived features did not
provide reliable
incremental predictive value over metadata for terminal ultimate load
under pooled evaluation. Moving to treatment- and campaign-level transfer
materially changed which conclusions the evidence could support, with
campaign holdout acting as a severe extrapolation stress test given how
closely campaign identity is entangled with specimen design here. The degradation and
proxy-remaining-useful-life work is methodological and diagnostic rather
than validated prognosis: no true remaining-useful-life supervision exists
in the dataset. A leakage-safe dataset for a subsequent four-class
severity classification task was also prepared, though not yet modelled.
The activity's principal contribution is a rigorously bounded account of
what repeated corrosion imagery can and cannot support here, together with
a traceable basis for the work that should follow.
